% =====================================================================
% USENIX-oriented source for:
%   Adversarial Intent Telemetry: An Empirical Study of
%   Behavioral-Signature Detection Under Adaptive Perturbation
%
% Derived (compressed) from the v9 single-column draft
% (../../Decentralized_Telemetry_Adversarial_AI_Intent_v8.1.tex), which
% remains the authoritative long-form draft. Content claims are identical;
% see paper/length_reduction_plan.md for what was cut and where it lives.
%
% The official USENIX style file (usenix.sty) is present in this directory
% and is used by default; the \IfFileExists switch below detects it. The
% official style already provides Times (mathptmx), T1 fontenc, utf8
% inputenc, microtype, cite, url/breakurl, xcolor, and colored hyperref, so
% those packages are NOT re-loaded when it is active (doing so clashes on
% hyperref options). The fallback branch (fired only if usenix.sty is
% removed) emulates USENIX geometry (10pt Times, two columns, 7in x 9.25in
% text block, 0.33in gutter) so lengths stay realistic during offline edits.
%
% Compile FROM THE REPOSITORY ROOT (paths to tables/figures are
% root-relative):
%   pdflatex -interaction=nonstopmode -halt-on-error \
%     -output-directory=paper/build -jobname=usenix_draft \
%     paper/usenix/adversarial_intent_telemetry_usenix.tex
%   bibtex paper/build/usenix_draft
%   pdflatex ... (twice more)
% =====================================================================
\documentclass[letterpaper,twocolumn,10pt]{article}

% Official USENIX style is loaded if found. Two detection paths so the same
% source builds both from paper/usenix/ (bare filename on the input path) and
% from the repository root (root-relative path; the documented compile CWD,
% where TeX does not search the main file's own directory). The style
% provides Times (mathptmx), fontenc, inputenc, microtype, cite, url/breakurl,
% xcolor, and colored hyperref, so those are not re-loaded here.
\IfFileExists{usenix.sty}{%
  \usepackage{usenix}%
}{%
  \IfFileExists{paper/usenix/usenix.sty}{%
    % \input (not \usepackage) because the file is off the input path; the
    % \makeatletter guard is required so usenix.sty's \if@twocolumn parses.
    \makeatletter\input{paper/usenix/usenix.sty}\makeatother
  }{%
    % ---- Fallback USENIX-geometry emulation (only if usenix.sty is absent) -
    \usepackage[letterpaper,total={7in,9.25in},centering]{geometry}
    \setlength{\columnsep}{0.33in}
    \usepackage{times}
    \usepackage[T1]{fontenc}
    \usepackage[utf8]{inputenc}
    \usepackage[hidelinks]{hyperref}
    \usepackage{xurl}
  }%
}

\usepackage{amsmath, amssymb}
\usepackage{booktabs}
\usepackage{makecell}
\usepackage{graphicx}
\usepackage{enumitem}
\setlist{itemsep=1pt, topsep=2pt, leftmargin=1.4em}
\renewcommand\theadfont{\bfseries\footnotesize}

% --- Anonymization switch (flip to \anonymoustrue for double-blind) --------
\newif\ifanonymous
\anonymousfalse

\title{\textbf{Adversarial Intent Telemetry: An Empirical Study of\\
Behavioral-Signature Detection Under Adaptive Perturbation}}
\ifanonymous
\author{Anonymous submission}
\else
\author{Fahrawn Gill\\ {\normalsize Advisor, AI Governance \& Cross-Platform
Safety, Alliance to Counter Crime Online (ACCO)}}
\fi
\date{}

\begin{document}
\maketitle

% =====================================================================
\begin{abstract}
Cross-platform safety infrastructure exchanges \emph{content} signatures at
production scale; a natural next step is to exchange \emph{behavioral}
signatures: structural fingerprints of adversarial conversations, shared
before harmful content exists. A prior design document specified such a
scheme: feature manifests matched across providers with banded MinHash
signatures, backed by a trajectory-level sequence model. This paper treats
that design as the object under test and asks whether it survives real
adversarial data and adaptive perturbation.

The answer is mixed, and the decomposition is the contribution. On the PAN
2012 corpus, a per-message linear classifier detects grooming conversations
well (AUC 0.986; recall 0.95 at FPR $\leq 0.05$). The design's core
cross-provider primitive does not: at the recommended banded-MinHash
operating point, recall on real conversations is 0.018 at FPR 0.0013, and
usable recall (0.92) appears only at FPR 0.99. The trajectory-level sequence
model does not beat the per-message baseline on clean data (F1 lift
$-0.023$, 95\% CI $[-0.042, -0.002]$). One robustness signal survives
stricter testing but changes attribution: under discourse-noise perturbation
sweeps the sequence model degrades less than the per-message baseline at
every strength, replicating under an author-disjoint split (heavy-noise AUC
drop 0.029 vs.\ 0.058) with paired bootstrap support, but matched controls
show the advantage is conversation-level averaging, not trajectory
structure: an order-invariant control is at least as robust, and order
shuffling leaves the sequence model unchanged. A false-positive substrate
analysis on structurally similar benign conversations found no elevated FPR
relative to random controls (0.030 vs.\ 0.055). Perturbation families are
rule-based, not adaptive, and PAN 2012 is a historical human-grooming
corpus, not evidence about agentic automation. The artifact includes
claim-consistency checks verifying every headline number against committed
result files.
\end{abstract}

% =====================================================================
\section{Introduction}
\label{sec:intro}

Content-hash exchange is the one piece of cross-platform trust-and-safety
infrastructure with an unambiguous production record: PhotoDNA-class
perceptual hashing sustains sub-millisecond matching at
ten-million-hash scale~\cite{photodna2023}, the Tech Coalition's Lantern
program has moved over two million safety signals across sixty member
companies, and the Video Hash Interoperability Project (VHIP) showed that
\emph{approximate}-match signatures interoperate across heterogeneous vendor
pipelines~\cite{lantern2026}. Meanwhile the abuse surface is moving upstream
of content~\cite{lantern2026,sok2026promptinjection,greshake2023}. It is
therefore tempting to extend the content-hash pattern one layer up: extract
structural features from a conversation, hash them into a compact signature,
and share it across providers so that adversarial \emph{behavior} detected
at one platform becomes detectable at its peers before content is generated.

That temptation should be tested, not assumed. Content hashes re-identify
near-identical objects. Behavioral signatures face a different problem: two
conversations executing the same adversarial strategy are \emph{different
texts}, and the question is whether any fixed structural representation
places them close together, and benign conversations far apart,
reliably enough to match at a deployable false-positive rate.

A prior design document (``Decentralized Telemetry for Adversarial AI
Intent,'' v8.1, included with the artifact~\cite{artifact2026}) specified a
complete scheme of this kind: canonical feature manifests, banded MinHash
signatures, a trajectory-level sequence model, trust-decay scheduling, and a
governance mapping. This paper does not announce that design; it evaluates
it. Several of its central claims do not survive. We state the mixed
result up front:

\begin{itemize}
\item A per-message baseline (LinearSVC over character n-grams) detects
grooming in PAN 2012 at AUC 0.986, recall 0.95 at FPR $\leq 0.05$:
behavioral structure in real adversarial conversations is detectable
(\S\ref{sec:rq1}).
\item The design's core cross-provider primitive fails on the same data:
banded MinHash matching at the recommended operating point recalls 1.8\% of
real predator conversations at FPR 0.0013; recall 0.92 requires FPR 0.99.
No usable operating point exists on the frontier (\S\ref{sec:rq2}). This is
a headline result, not a footnote.
\item The trajectory-level sequence model does not improve on the
per-message baseline on clean data: F1 lift $-0.023$, 95\% CI
$[-0.042, -0.002]$ (\S\ref{sec:rq3}).
\item The one positive result, with its attribution corrected: under
discourse-noise perturbation sweeps the sequence model's AUC drop is smaller
than the baseline's at every strength, replicating author-disjoint
(0.029 vs.\ 0.058 at heavy noise) with paired bootstrap support. Matched
aggregation and order controls then show the advantage comes from
\emph{conversation-level averaging}, not trajectory or order information:
an order-invariant twin of the sequence model is at least as robust, and
order shuffling barely moves the sequence model (\S\ref{sec:rq4}).
\item A false-positive substrate analysis on structurally similar benign
conversations, the population a deployed system would most plausibly harm,
found no elevated FPR relative to random benign controls on the PAN
proxy (\S\ref{sec:rq5}).
\end{itemize}

\textbf{Contributions.}
(1) An evaluation of a proposed behavioral-signature telemetry design under
real-data and perturbation tests on PAN 2012.
(2) A negative result for the banded-MinHash signature primitive: no usable
recall/FPR operating point on real text.
(3) A negative result for clean-data trajectory lift over a strong
per-message baseline.
(4) A perturbation-robustness result that survives author-disjoint
replication with paired statistical support, and matched aggregation/order
controls showing it is conversation-level averaging, not trajectory
structure, correcting the design's trajectory hypothesis.
(5) A substrate-conditional false-positive measurement on structurally
similar benign conversations, with no elevated FPR observed on the proxy.
(6) A reproducible artifact with claim-consistency checks, machine-generated
result tables, provenance documentation, and a validity-boundary
statement~\cite{artifact2026}.

Table~\ref{tab:evidence} maps every claim to its evidence status and source
file; nothing in this paper should be read as stronger than its row there.

\begin{table*}[!t]
\centering
\footnotesize
\renewcommand{\arraystretch}{1.15}
\begin{tabular}{@{}p{5.6cm} p{3.6cm} p{7.6cm}@{}}
\toprule
\thead{Claim} & \thead{Status} & \thead{Evidence (result file, section)} \\
\midrule
Per-message classifier detects grooming on clean PAN 2012 & Real data: holds & AUC 0.986, recall 0.95 @ FPR $\leq 0.05$; conversation-level split. \texttt{trajectory\_lift.json} (\S\ref{sec:rq1}) \\
MinHash matching separates real conversations at deployable FPR & Real data: \textbf{negative} & Recall 0.018 @ FPR 0.0013 at $(16,16)$; recall 0.92 only @ FPR 0.99; author-disjoint. \texttt{m3\_frontier.json} (\S\ref{sec:rq2}) \\
Trajectory model lifts F1 over baseline on clean data & Real data: negative/inconclusive & Lift $-0.023$, CI $[-0.042,-0.002]$. \texttt{trajectory\_lift.json} (\S\ref{sec:rq3}) \\
Sequence model shows a smaller AUC drop than per-message baseline & Real data: holds (rule-based families) & Smaller AUC drop at all strengths, both splits, 5 seeds; paired CIs exclude zero. \texttt{perturbation\_sweep*.json}, \texttt{perturbation\_paired\_tests.json} (\S\ref{sec:rq4}) \\
The robustness reflects trajectory/order information & Real data: \textbf{negative} & Order-invariant control at least as robust; shuffled-order control tracks the model. Attribution: conversation-level averaging. \texttt{perturbation\_model\_controls.json} (\S\ref{sec:rq4}) \\
Structurally similar benign conversations absorb disproportionate FPs & Real data (proxy): not observed & Classifier FPR 0.030 vs.\ 0.055 on random controls; author-disjoint, 3 seeds. \texttt{fp\_substrate.json} (\S\ref{sec:rq5}) \\
Banded-MinHash S-curve follows the closed form & Synthetic: sanity check only & Implementation matches Eq.~\eqref{eq:lsh-banded}; no real-data utility claim (\S\ref{sec:design}) \\
Byzantine tolerance; SPRT vs.\ Hoeffding isolation & Simulation only & $\beta^\star{=}0.5$ (stealth 0.4); $2.27\times$ speedup with a weak comparator caveat. \texttt{m8\_*.json} (\S\ref{sec:rq6}) \\
Payoff-perturbation mechanism raises cooperation & Analytical only & Game matrices, not LLM behavior. \texttt{f3\_reciprocity.json} (\S\ref{sec:rq6}) \\
A2 agentic automation as the operative adversary & Design target: not validated & Motivated by reporting~\cite{lantern2026}; data are 2012 human conversations (\S\ref{sec:threat}) \\
Manifest schema, handshake, latency/volumetric targets & Design context: proposed, not validated & Retained from the design document; no implementation evaluated (\S\ref{sec:design}) \\
\bottomrule
\end{tabular}
\caption{Evidence and validity matrix. \emph{Simulation} and \emph{analytical} rows exercise mechanisms under stated parameters, not deployments; \emph{design context} rows carry no empirical claim.}
\label{tab:evidence}
\end{table*}

% =====================================================================
\section{Background and Motivation}
\label{sec:background}

\textbf{Content-hash exchange works in production.} Lantern is the operating
cross-platform child-safety signal exchange (sixty members, two million
signals, 350{,}000 enforcement actions by end of 2025, externally reviewed
governance~\cite{lantern2026,bsrhria}); VHIP extends the pattern to
approximate matching across heterogeneous video-hash pipelines (435{,}000
videos in 2025)~\cite{lantern2026}. VHIP is the closest architectural
precedent for the design evaluated here, and the analogy it suggests is
exactly what this paper tests at the behavioral layer.

\textbf{The abuse surface is moving upstream of content.} The 2025--2026
reporting is directionally clear: two-orders-of-magnitude growth in
AI-generated abuse video, over one million generative-AI-related CyberTipline
reports in nine months, seventeen-fold growth in AI-related Lantern
signals~\cite{lantern2026}; in parallel, the agentic-security literature
documents indirect prompt injection~\cite{greshake2023}, cross-tool
poisoning across MCP clients~\cite{mcpclients2026}, zero-click
exfiltration~\cite{echoleak2025}, and protocol-level trust-model
failures~\cite{sok2026promptinjection,mcprce2026,breakingprotocol2026}.
% TODO(bib): replace mcprce2026 (press) with vendor advisories / NVD entries.
Two boundaries on this motivation: the figures motivate the \emph{design
goal} and validate nothing; and the hypothesized agentic adversary is not
the adversary in our data: PAN 2012 records human grooming from 2012, and
where 2025 NCMEC figures appear at all they are sampling priors for
perturbation construction, never evidence of an agentic-abuse distribution.

\textbf{Why behavioral signatures are harder than content hashes.} A
perceptual hash re-identifies a near-duplicate object. A behavioral
signature must place two \emph{different} conversations sharing a strategy
close together while keeping the enormous benign population, including
benign conversations about ages, meetings, secrets, and relationships,
far away. Whether any fixed manifest-and-hash representation achieves that
on real text is an empirical question; \S\ref{sec:rq2} answers it negatively
for the evaluated design.

% =====================================================================
\section{Threat Model and Design Under Test}
\label{sec:threat}

\subsection{Threat Model}
The threat model is retained from the design under test, with one framing
correction made explicit: the adversary classes define what the design
\emph{targets}, not what this paper \emph{validates}. The protected assets
are end-user safety (priority to minors), platform policy-enforcement
integrity, regulatory standing, and signature unlinkability. Five adversary
classes: \textbf{A1}, an individual offender using structural prompt framing
against a hosted model; \textbf{A2}, an agentic automation network managing
adversarial conversations at scale across platforms with persistent persona,
the design's lead class, motivated by reporting~\cite{lantern2026} but
\emph{not validated in this study}, whose data are 2012 human conversations;
\textbf{A3}, a compromised or strategic federation participant emitting
signatures that do not reflect genuine detections; \textbf{A4}, an external
observer attempting dictionary attacks on emitted signatures; \textbf{A5}, a
regulatory misuser invoking the scheme outside its stated scope. The design
states five goals (federation detection lift; unlinkability; Byzantine
tolerance; schema-level data minimization; scope confinement); this paper
measures primitives relevant to detection lift and its false-positive side
(\S\ref{sec:rq1}--\ref{sec:rq5}) and Byzantine tolerance in simulation only
(\S\ref{sec:rq6}). Out of scope, unchanged from the design and central
to the ethics posture (\S\ref{sec:ethics}): end-to-end encrypted content
(no client-side scanning is proposed), fully airgapped open-source
deployments, and harms outside the enumerated scope.

\subsection{Design Under Test}
\label{sec:design}
The design's intuition: adversarial conversations have detectable
\emph{form} that survives lexical evasion; that form is captured as a
canonical feature manifest (inferred intent class, behavioral phase, anomaly
bins, role-injection markers; full schema in the artifact's
\texttt{spec/manifest-schema.json} and the design document); manifests are
matched approximately across providers with banded MinHash; and
trajectory-level structure is the layer an adversary cannot cheaply rewrite
away. The design encodes detections into six structural classes (semantic
reframing, lexical substitution, compositional decomposition, authority/role
injection, hypothetical sandboxing, cumulative context poisoning); in this
paper that taxonomy is design vocabulary only, not a validated
classification~\cite{medredteam2026}. The design is a runtime layer,
orthogonal to training-time
alignment~\cite{opencharacter2025,inoculation2025}.

\textbf{Signature primitive.} Cryptographic hashes have zero tolerance for
input drift across heterogeneous providers, so the design adopts
locality-sensitive hashing~\cite{broder1997}. With $J$ the Jaccard
similarity between two feature manifests and a MinHash signature of length
$L = b\cdot r$ split into $b$ bands of $r$ rows, the candidate-match
probability is the S-curve
\begin{equation}
P_{\text{match}}(J;\, b, r) \;=\; 1 - (1 - J^{r})^{b},
\label{eq:lsh-banded}
\end{equation}
with inflection $J^\star \approx (1/b)^{1/r}$; the per-comparison FPR is
bounded by $P_{\text{match}}$ at the maximum benign Jaccard. The design
recommended $(b,r) = (16,16)$ at $L = 256$, placing $J^\star \approx 0.841$.
The unstated empirical premise, that manifests of related adversarial
conversations achieve Jaccard near $J^\star$ while benign pairs stay below,
is what \S\ref{sec:rq2} tests. A synthetic sanity check confirms the
implementation follows Eq.~\eqref{eq:lsh-banded} (within 0.05 at every
tested $J$) and already hints at the problem: even under synthetic Jaccard
distributions constructed to be favorable, recall at the recommended point
is only 0.085. The design's emission/receive/match protocol, wire format,
decision rule (a match enqueues independent human investigation, never
automatic enforcement), and privacy invariant (no raw text or user identity
in any payload) are retained in the design document and the long-form
draft; they are design context, not evaluated systems.

\textbf{Trajectory hypothesis.} Grooming escalates through documented
phases; the design asserted that phase \emph{transitions} carry the
detection signal and are structurally expensive to evade (destroying the
trajectory destroys the attack), treating that layer as the one an adversary
cannot cheaply rewrite away.
Empirically this hypothesis earns: no clean-data lift (\S\ref{sec:rq3}), a
small order-sensitivity asymmetry, and a robustness advantage that matched
controls attribute to conversation-level averaging rather than order
information (\S\ref{sec:rq4}).

% =====================================================================
\section{Data and Methods}
\label{sec:methods}

All numbers are produced by scripts in the artifact's \texttt{experiments/}
directory and written to JSON result files; each result table names its
source file, tables are machine-generated from the JSONs, and
\texttt{scripts/check\_claims.py} verifies the headline numbers
(\S\ref{sec:openscience}).

\textbf{Data.} The PAN 2012 Sexual Predator Identification
corpus~\cite{villatoro2012} (66{,}927 chat conversations from the training
partition, roughly 3\% involving a labeled predator author) is the
substrate for all real-data claims. PAN 2012 is public, open research data
and also sensitive online-safety material: the artifact does not
redistribute it, no raw dialogue appears in this paper or any output, and
all reported quantities are aggregate metrics, structural labels, or counts.
GT-HarmBench~\cite{gtharmbench2026}, a gated dataset of game-theoretic
scenarios, supplies payoff matrices for the analytical F3 result only.
% TODO(bib): verify gtharmbench2026 identifier; load-bearing for S5.6.

\textbf{Splits.} The \emph{author-disjoint} 80/20 split assigns all of a
predator author's conversations to one side (113 train / 29 test authors;
13{,}202 test conversations, 219 positive; benign conversations, which have
no predator author, split 80/20 by conversation ID) and is used by the
MinHash frontier, the FP-substrate analysis, the sweep replication, and all
control experiments. The \emph{stratified conversation-level} split (13{,}387
test conversations, 404 positive) is used by the trajectory-lift experiment
and the original sweep and is flagged wherever used. Split membership is
fixed across perturbation conditions (split-ID hashes and a
zero-author-overlap check in the sanity files).

\textbf{Models.} Per-message baseline: TF-IDF over character 3--5-grams
(10{,}000 features), LinearSVC on message labels, conversation scored by the
mean message decision value. Sequence model: logistic regression over a
position-weighted mean of message TF-IDF vectors, a deliberately simple
instantiation of the trajectory idea. Controls (\S\ref{sec:rq4}): max and
top-5-mean aggregations of the same message scores; an order-invariant
concatenated-text classifier; an order-invariant \emph{conversation-mean}
twin of the sequence model (identical except position weights removed, so
their comparison isolates order information); and the sequence model scored
on order-shuffled messages. Per-model thresholds are calibrated once on the
clean test set at FPR $\leq 0.05$ and held fixed.

\textbf{MinHash frontier.} Conversations are shingled into $L=256$ MinHash
signatures; the 1{,}720 valid training-positive signatures form the
shared reference pool; each test conversation queries the pool with banded
matching at five $(b,r)$ settings. This instantiates the design's
cross-provider match primitive on real text.

\textbf{Perturbation sweeps.} Models are trained once on clean data; test
positives are perturbed at inference (an adaptive operator perturbs their
own conversations; negatives stay clean) at strengths
$p \in \{0.15, 0.35, 0.60\}$, five seeds each. Family 1, \emph{discourse
noise}: message truncation, retrieval-swap of training-split benign
messages, message drop (with a floor guard), adjacent-transposition order
jitter. Family 2, \emph{surface rewrite} (rule-based, local): bidirectional
chat-register substitutions from a fixed neutral table, casing/punctuation
jitter, message split/merge. Sanity files record per-condition modification
rates, token-Jaccard, order changes, and length shifts.

\textbf{Paired statistics.} For each strength and comparator, a stratified
paired bootstrap ($B{=}1000$; positives/negatives resampled separately with
identical indices for every model; perturbed scores seed-averaged) yields
95\% CIs for (sequence AUC drop $-$ comparator AUC drop); a CI below zero
favors the sequence model.

\textbf{FP substrate.} From author-disjoint test negatives, a
\emph{structurally similar benign} pool (3{,}979 conversations) is selected
by neutral lexical filters (age discussion, meeting plans,
secrecy/isolation, romantic language), never classifier predictions,
against a random benign control pool (9{,}004); three seeds sample 500 per
set; both detectors are evaluated with Wilson 95\% intervals. PAN negatives
are a proxy for sensitive benign populations; the corpus has no population
labels.

\textbf{Simulation and analytical components.} Reputation-weighted
aggregation is simulated under coordinated mis-emission with SPRT or
EMA-plus-Hoeffding isolation (parameters in the result files); the F3
mechanism is applied to GT-HarmBench $2\times2$ matrices analytically. No
LLM behavior is measured in either.

% =====================================================================
\section{Results}
\label{sec:results}

\subsection{RQ1: Per-message detection on clean data}
\label{sec:rq1}
Strong. The per-message LinearSVC reaches AUC 0.986 with recall 0.948 at
FPR $\leq 0.05$ (Table~\ref{tab:clean}). Behavioral structure in real
adversarial conversations is detectable at the local-classifier level. Two
qualifications: precision at that operating point is 0.412 (most alerts
are false positives at PAN's $\sim$3\% base rate), and nothing here speaks
to robustness (\S\ref{sec:rq4}) or cross-provider matching
(\S\ref{sec:rq2}).

\begin{table}[!t]
\centering
\footnotesize
\resizebox{\columnwidth}{!}{%
% Auto-generated by scripts/build_paper_tables.py
% Source of truth: experiments/results/trajectory_lift.json
% Do not edit by hand; re-run the script instead.
\begin{tabular}{@{}l c c c c@{}}
\toprule
Model & AUC & Recall & Precision & F1 \\
\midrule
Per-message LinearSVC & 0.986 & 0.948 & 0.412 & 0.575 \\
Sequence model & 0.977 & 0.933 & 0.392 & 0.552 \\
\midrule
\multicolumn{5}{@{}l}{Sequence-over-baseline F1 lift: $-0.023$, bootstrap 95\% CI $[-0.042,\,-0.002]$} \\
\bottomrule
\end{tabular}
}
\caption{Clean-data detection (conversation-level split; operating points at
FPR $\leq 0.05$). Source: \texttt{trajectory\_lift.json}.}
\label{tab:clean}
\end{table}

\subsection{RQ2: MinHash matching on real conversations}
\label{sec:rq2}
Fails; this is a headline finding. Table~\ref{tab:frontier} and
Figure~\ref{fig:frontier} give the full operating frontier on the
author-disjoint split (13{,}202 test conversations against the pool of
1{,}720 training-positive signatures). At the design-recommended $(16,16)$,
recall is \textbf{0.018} at FPR 0.0013; recall does not become usable until
$(128,2)$, where it is 0.92 at FPR \textbf{0.988}: the detector fires on
essentially everything. No point combines usable recall with tolerable FPR.
The mechanism is visible in the Jaccard statistics: real conversations
executing related adversarial behavior do not share enough surface
structure; inter-conversation Jaccard sits far below the $J^\star \approx
0.841$ inflection, so the S-curve, which behaves exactly as designed,
is evaluated in a region where it correctly returns no-match. The primitive
is not broken; the design's assumption about where real adversarial Jaccard
mass lies is falsified. A successor must change the representation or
abandon set-similarity matching for this layer.

\begin{table}[!t]
\centering
\footnotesize
% Auto-generated by scripts/build_paper_tables.py
% Source of truth: experiments/results/m3_frontier.json
% Do not edit by hand; re-run the script instead.
\begin{tabular}{@{}c c c c c@{}}
\toprule
$b$ & $r$ & Inflection $J^\star$ & Recall & FPR \\
\midrule
8 & 32 & 0.937 & 0.0046 & 0.0012 \\
16 & 16 & 0.841 & 0.0183 & 0.0013 \\
32 & 8 & 0.648 & 0.2557 & 0.0491 \\
64 & 4 & 0.354 & 0.7169 & 0.3294 \\
128 & 2 & 0.088 & 0.9224 & 0.9880 \\
\bottomrule
\end{tabular}

\caption{Banded-MinHash frontier on real PAN 2012 (author-disjoint split,
$L=256$). Source: \texttt{m3\_frontier.json}.}
\label{tab:frontier}
\end{table}

\begin{figure}[!t]
\centering
% Column-width figure generated from m3_frontier.json by
% scripts/build_paper_figures.py (deterministic; no experiment re-run).
\includegraphics[width=\columnwidth]{paper/figures/fig_frontier_column.pdf}
\caption{Recall/FPR frontier of the signature primitive on real text
(symlog FPR axis; points labeled with $(b,r)$). Generated from
\texttt{m3\_frontier.json}.}
\label{fig:frontier}
\end{figure}

\subsection{RQ3: Clean-data trajectory lift}
\label{sec:rq3}
Negative. The sequence model reaches AUC 0.977 vs.\ the baseline's 0.986;
F1 lift at the FPR $\leq 0.05$ operating point is $-0.023$ with bootstrap
95\% CI $[-0.042, -0.002]$ (Table~\ref{tab:clean}): the interval lies
below zero. This falsifies \emph{this} position-weighted representation as a
clean-data improvement, not trajectory modeling in general. One structural
asymmetry survives: the baseline is order-invariant by construction, while
the sequence model loses about one point of recall (0.933 $\to$ 0.923) under
adversarial order shuffling, and an adversary who destroys conversational
order also destroys the escalation structure the attack requires. That
asymmetry is what the perturbation experiments probe next.

\subsection{RQ4: Robustness under perturbation is real, but it is
aggregation, not trajectory}
\label{sec:rq4}

\textbf{The sweep and its author-disjoint replication.}
Table~\ref{tab:sweeps} reports both splits. On the conversation-level split
the baseline's AUC falls 0.986 $\to$ 0.912 at heavy noise while the sequence
model falls 0.977 $\to$ 0.940; the sequence drop is smaller at every
strength (heavy: 0.037 vs.\ 0.075). Under the author-disjoint split, clean
AUCs compress (0.975/0.971) and the ordering is unchanged (heavy drop 0.029
vs.\ 0.058, roughly half). The result was not author leakage. Earlier
single-point observations (discourse noise 0.92 $\to$ 0.79; an
NCMEC-prior-constrained set collapsing the per-message signal to 0.13 while
the sequence model retained 0.57) are consistent but superseded by the
dose-response sweeps.

\begin{table}[!t]
\centering
\footnotesize
\resizebox{\columnwidth}{!}{%
% Auto-generated by scripts/build_paper_tables.py
% Source of truth: experiments/results/perturbation_sweep.json + perturbation_sweep_author_disjoint.json
% Do not edit by hand; re-run the script instead.
\begin{tabular}{@{}l cc cc@{}}
\toprule
 & \multicolumn{2}{c}{Conversation-level split} & \multicolumn{2}{c}{Author-disjoint split} \\
\cmidrule(lr){2-3}\cmidrule(lr){4-5}
Strength & Per-msg & Sequence & Per-msg & Sequence \\
\midrule
none & 0.986 & 0.977 & 0.975 & 0.971 \\
light & 0.971 (0.016) & 0.968 (0.010) & 0.963 (0.012) & 0.964 (0.006) \\
medium & 0.949 (0.037) & 0.955 (0.022) & 0.947 (0.028) & 0.958 (0.013) \\
heavy & 0.912 (0.075) & 0.940 (0.037) & 0.917 (0.058) & 0.941 (0.029) \\
\bottomrule
\end{tabular}
}
\caption{Discourse-noise sweeps on both splits: mean AUC over 5 seeds (AUC
drop from clean in parentheses). Sources: \texttt{perturbation\_sweep.json},
\texttt{perturbation\_sweep\_author\_disjoint.json}.}
\label{tab:sweeps}
\end{table}

\textbf{Matched controls change the attribution.} Table~\ref{tab:controls}
evaluates seven models/controls on the same author-disjoint split, perturbed
conversations, and seeds; Table~\ref{tab:paired} gives paired bootstrap CIs
on the AUC-drop differences. Two facts. First, the sequence model's
advantage is statistically resolved against every per-message comparison:
mean, max, top-5, and the concatenation classifier (the best clean model at
AUC 0.989, but the \emph{worst} degrader at heavy noise, 0.087), with all
CIs excluding zero in its favor. Second, the order-invariant
conversation-mean control, identical to the sequence model except that
position weights are removed, degrades \emph{less} than the sequence model
at every strength (paired CI excludes zero in the control's favor; heavy:
$+0.004$, CI $[+0.002, +0.006]$), and the shuffled-order control tracks the
sequence model almost exactly (heavy drop 0.030 vs.\ 0.029). The tested
sequence model barely uses order, and its position weights cost a small
amount of robustness rather than buying any.

\begin{table}[!t]
\centering
\footnotesize
\resizebox{\columnwidth}{!}{%
% Auto-generated by scripts/build_paper_tables.py
% Source of truth: experiments/results/perturbation_model_controls.json
% Do not edit by hand; re-run the script instead.
\begin{tabular}{@{}l c c c c@{}}
\toprule
 & Clean & \multicolumn{3}{c}{AUC drop at strength} \\
\cmidrule(lr){3-5}
Model / control & AUC & light & medium & heavy \\
\midrule
Per-message mean (baseline) & 0.975 & 0.012 & 0.028 & 0.058 \\
Per-message max & 0.895 & 0.008 & 0.016 & 0.036 \\
Per-message top-5 mean & 0.934 & 0.011 & 0.026 & 0.051 \\
Conversation concat (order-inv.) & 0.989 & 0.010 & 0.024 & 0.087 \\
Conversation mean (order-inv.) & 0.968 & 0.005 & 0.010 & 0.026 \\
\textbf{Sequence (position-weighted)} & 0.971 & 0.006 & 0.013 & 0.029 \\
Sequence, shuffled order (control) & 0.972 & 0.006 & 0.012 & 0.030 \\
\bottomrule
\end{tabular}
}
\caption{Aggregation/order controls (author-disjoint, discourse noise, 5
seeds): clean AUC and mean AUC drop per strength. Source:
\texttt{perturbation\_model\_controls.json}.}
\label{tab:controls}
\end{table}

\begin{table}[!t]
\centering
\scriptsize
\resizebox{\columnwidth}{!}{%
% Auto-generated by scripts/build_paper_tables.py
% Source of truth: experiments/results/perturbation_paired_tests.json
% Do not edit by hand; re-run the script instead.
\begin{tabular}{@{}l c c c@{}}
\toprule
Comparison & light & medium & heavy \\
\midrule
vs.\ per-message mean (baseline) & $[-0.0202, -0.0054]$$^*$ & $[-0.0239, -0.0080]$$^*$ & $[-0.0365, -0.0165]$$^*$ \\
vs.\ per-message max & $[-0.0217, -0.0052]$$^*$ & $[-0.0198, -0.0043]$$^*$ & $[-0.0471, -0.0193]$$^*$ \\
vs.\ per-message top-5 mean & $[-0.0213, -0.0073]$$^*$ & $[-0.0290, -0.0141]$$^*$ & $[-0.0483, -0.0263]$$^*$ \\
vs.\ conversation concat (order-inv.) & $[-0.0138, -0.0036]$$^*$ & $[-0.0217, -0.0083]$$^*$ & $[-0.0657, -0.0448]$$^*$ \\
vs.\ conversation mean (order-inv.) & $[+0.0004, +0.0034]$$^*$ & $[+0.0007, +0.0050]$$^*$ & $[+0.0018, +0.0064]$$^*$ \\
\bottomrule
\end{tabular}
}
\caption{Paired stratified bootstrap ($B{=}1000$) 95\% CIs for (sequence AUC
drop $-$ comparator drop); negative favors the sequence model; $^*$ = CI
excludes zero. Source: \texttt{perturbation\_paired\_tests.json}.}
\label{tab:paired}
\end{table}

\begin{figure}[!t]
\centering
% Column-width figure generated from perturbation_model_controls.json by
% scripts/build_paper_figures.py (deterministic; no experiment re-run).
\includegraphics[width=\columnwidth]{paper/figures/fig_controls_column.pdf}
\caption{AUC vs.\ perturbation strength for all models/controls
(author-disjoint split). Generated from
\texttt{perturbation\_model\_controls.json}.}
\label{fig:controls}
\end{figure}

\textbf{A second perturbation family does not discriminate.} Under the
rule-based surface-rewrite/segmentation family, every model is essentially
unaffected (largest mean AUC drop under 0.002;
\texttt{perturbation\_second\_family.json}). Character n-gram features are
insensitive to this rewriting: mild good news for all detectors, but no
evidence either way on the robustness ordering, which therefore rests on the
discourse-noise family.

\textbf{What RQ4 supports.} \emph{Conversation-level averaging, training
and scoring on features averaged over the whole conversation, degrades
roughly half as fast as per-message scoring under discourse noise, on both
splits, with paired statistical support; the order information in the tested
sequence model contributes nothing to that robustness.} The families are
rule-based, not adaptive; paraphrase attacks are untested; and a genuinely
order-sensitive architecture remains unevaluated: the trajectory
hypothesis proper is still open, but it has no support from this data.

\subsection{RQ5: False positives on structurally similar benign
conversations}
\label{sec:rq5}
Not elevated on this proxy: the effect runs the other way. Benign
conversations that surface-resemble grooming transcripts (age talk, meeting
plans, secrecy, romantic language) proxy the populations a deployed detector
would most plausibly harm. Table~\ref{tab:fps}: the classifier's FPR on the
keyword-matched set is \textbf{0.030} vs.\ \textbf{0.055} on random controls,
consistently across seeds; the MinHash primitive flags almost nothing in
either set (0 and 1 conversations total), as RQ2 predicts. This weakens the
assumption that surface-similar benign traffic is automatically higher-risk
under a lexical detector (possibly because keyword-matched negatives
resemble the classifier's training negatives), but it does not clear the
deployment concern: PAN negatives cannot label the real populations of
concern, and external sensitive-population substrates (under ethics review)
would be more probative. What it establishes is the method:
substrate-conditional FPR measurement with controls is cheap, and any
deployment claim would need it on real substrates.

\begin{table}[!t]
\centering
\footnotesize
\resizebox{\columnwidth}{!}{%
% Auto-generated by scripts/build_paper_tables.py
% Source of truth: experiments/results/fp_substrate.json
% Do not edit by hand; re-run the script instead.
\begin{tabular}{@{}l c c@{}}
\toprule
Detector & \makecell{Structurally similar benign \\ FPR [pooled Wilson 95\% CI]} & \makecell{Random benign control \\ FPR [pooled Wilson 95\% CI]} \\
\midrule
MinHash $(b{=}16, r{=}16)$ & 0.000 [0.000,\,0.003] & 0.001 [0.000,\,0.004] \\
LinearSVC @ FPR$\leq$0.05 & 0.030 [0.022,\,0.040] & 0.055 [0.045,\,0.068] \\
\bottomrule
\end{tabular}
}
\caption{FP substrate analysis (author-disjoint; $n{=}500$/set; 3 seeds;
pooled Wilson 95\% CIs). Source: \texttt{fp\_substrate.json}.}
\label{tab:fps}
\end{table}

\subsection{RQ6: Simulation and analytical components}
\label{sec:rq6}
These exercise mechanisms under stated parameters; they are not deployment
evidence and must not be aggregated with the real-data findings.
\emph{Byzantine tolerance (simulation):} network detection quality degrades
$<5\%$ up to an empirical Byzantine fraction $\beta^\star = 0.5$ (stealth
adversary: 0.4). \emph{SPRT vs.\ Hoeffding (simulation, weak comparator):}
the sequential test isolates mis-emitting senders $2.27\times$ faster (mean
isolation step) with SPRT false-isolation $\leq 0.044$, but the Hoeffding
baseline's own false-isolation rate is extreme ($\approx 0.996$) because it
re-applies a fixed-$n$ threshold to a high-variance EMA at every step with
no exoneration state; this is expected from the baseline's definition, the
IID fixed-$n$ assumption behind the bound does not hold sequentially, and
the false-isolation gap should not be quoted as a finding.
\emph{F3 reciprocity (analytical):} the design's payoff-perturbation
mechanism raises equilibrium cooperation on
GT-HarmBench~\cite{gtharmbench2026} $2\times2$ matrices by 0.184 overall and
cuts Prisoner's Dilemma defection from 1.0 to 0.57. Mechanism design on
matrices~\cite{ostrom1990}; nothing about LLM or provider behavior.

% =====================================================================
\section{Discussion}
\label{sec:discussion}

\textbf{What the decomposition means.} The design bundled three bets:
(i) behavioral structure is detectable; (ii) it can be shared as approximate
set-similarity signatures; (iii) trajectory modeling is the robust layer.
Bet (i) holds on clean historical data (RQ1). Bet (ii) fails on real text
(RQ2), at the representation level, not the primitive level, which
redirects rather than ends the work: the open problem is a representation
under which related-but-distinct adversarial conversations achieve high
similarity while benign traffic does not, and RQ5 shows benign structural
look-alikes are not automatically the binding constraint. Bet (iii) fails in
the form the design stated it: no clean lift (RQ3), and the robustness that
appeared to support it is attributable to conversation-level averaging
(RQ4). What survives is smaller but useful to a detection engineer: train
and score at the conversation level over averaged features and the detector
degrades roughly half as fast under discourse noise. Whether a genuinely
order-sensitive architecture can add robustness \emph{beyond} that
aggregation effect is now the precise, testable form of the trajectory
hypothesis.

\textbf{False positives frame everything.} At PAN's class balance even the
strong RQ1 classifier yields precision 0.41 at 5\% FPR: most alerts are
false. Any system built on these primitives is a review-queue prioritizer,
not an enforcement trigger; the design's own rule that a match enqueues
independent human investigation is the minimum viable posture. The design's
failure-mode catalog (signature pollution, FP cascades, free-riding, scope
creep, replay/dictionary attacks, trust-root failure) motivated the M8
simulations and the RQ5 method; its mitigations (trust decay, per-class
FP ceilings, TTLs, rate limits, minimum-entropy gating, multi-root PKI)
remain design context, with one measured correction: the Hoeffding-derived
isolation rule false-isolates almost every honest sender when applied
sequentially (\S\ref{sec:rq6}), so any deployment of it would need
sequential correction (which is what the SPRT variant is).

% =====================================================================
\section{Limitations}
\label{sec:limitations}
(1) \emph{Corpus vintage and adversary mismatch.} PAN 2012 records human
grooming in 2012-era chat; transfer to current platforms, other languages,
and hypothesized agentic (A2) operations is untested. (2) \emph{Rule-based
perturbation families only.} The robustness ordering rests on the
discourse-noise family; the surface-rewrite family left every model
unaffected and is uninformative for the comparison; paraphrase/LLM-rewrite
and genuinely adaptive attacks are untested and could reverse the ordering.
(3) \emph{Split coverage.} The robustness result replicates author-disjoint,
but clean-data RQ1/RQ3 numbers are reported on the conversation-level split
and may be optimistic (author-disjoint clean AUCs are lower: 0.975/0.971).
(4) \emph{Weak sequence model, with a measured consequence.} The tested
representation barely uses order (shuffled control tracks it; position
weights slightly hurt robustness), so RQ4's attribution finding is about
this model class; a genuinely order-sensitive encoder is unevaluated in both
directions. (5) \emph{Proxy substrate.} RQ5's benign substrate is
keyword-selected PAN negatives, not real sensitive populations, and seeds
resample overlapping pools. (6) \emph{Simulation/analytical parameters.}
M8 and F3 hold under stated parameters only.

% =====================================================================
\section{Conclusion}
\label{sec:conclusion}
We evaluated a proposed cross-provider behavioral-signature telemetry design
against real adversarial data and adaptive perturbation. The design's core
cross-provider primitive fails on real conversations at any deployable
false-positive rate; its trajectory hypothesis earns no support in either
tested form: no clean-data lift, and a perturbation-robustness advantage
that matched controls attribute to conversation-level averaging rather than
order information. What survives is a decomposition useful to anyone
building in this space: per-message behavioral detection works on clean
historical data; conversation-level averaging is the robust way to score it
under discourse noise; benign structural look-alikes did not absorb
disproportionate false positives on the available proxy; and every negative
result is documented with the same care as the positive ones. The evidence
supports no deployment claim. Two open questions are now precise: a
representation under which real adversarial conversations actually cluster,
and an order-sensitive architecture that beats the conversation-mean control
under the same paired-test framework. We state both so that this design, or
its successor, can be falsified faster next time.

% =====================================================================
\section{Ethical Considerations}
\label{sec:ethics}

\textbf{Data handling.} All real-data experiments use PAN 2012, which is
public, open research data that is also sensitive online-safety material
involving real victims and offenders. The artifact does not redistribute the corpus;
no raw dialogue appears in this paper, the artifact's documentation, or any
committed result file; all reported quantities are aggregate metrics,
structural labels, and counts. No abuse content, grooming scripts, or
operational evasion material was generated: the perturbation families are
generic discourse/surface noise, not offender tradecraft, and no corpus text
was sent to external services.

\textbf{No client-side scanning; no E2EE weakening.} Neither the design
under test nor this evaluation involves scanning end-to-end encrypted
content or client devices, and nothing here proposes weakening encryption;
everything operates over data a platform can already observe under its own
policies.

\textbf{No deployment recommendation.} The evidence would not support one:
the signature primitive fails on real data, and the robustness finding
spans rule-based families only. The design document's latency, throughput,
and topology targets are design assumptions from production
analogues~\cite{photodna2023}; no implementation was built or measured, and
none of those numbers is a claim of this paper.

\textbf{False-positive harms are the central deployment risk.} Behavioral
detectors threaten specific benign populations whose conversations share
surface structure with abuse transcripts: minors discussing age and
relationships, LGBTQ+ youth seeking support, mental-health peer support,
harm-reduction discussion, journalism, and legitimate research.
\S\ref{sec:rq5} measures a proxy and finds no elevated FPR, but the proxy
cannot label the real populations, and nothing here clears the risk; any
deployment-oriented study would need substrate-conditional FP measurement on
real sensitive-population data under ethics review.

\textbf{Human review, appeal, no automated enforcement.} A signature or
classifier match must never trigger enforcement automatically; the design's
own enqueue-for-independent-human-investigation rule is the minimum posture,
fully automated enforcement would additionally be in tension with GDPR
Article~22, and any future system would require an accessible appeals path
and multi-stakeholder scope governance with civil-society representation on
the model of Lantern's externally reviewed governance~\cite{bsrhria}.

\textbf{No compliance claims.} This work does not establish compliance with
the EU AI Act, GDPR, the UK Online Safety Act, or any other
regime~\cite{euaiact2026,gpaicop2026}; conformity assessment is
participant-level work. The design document's compliance mapping is
superseded by these non-claims.

% =====================================================================
\section{Open Science}
\label{sec:openscience}

\ifanonymous
The artifact is available to reviewers via the anonymized link provided in
the submission system. % TODO(anonymization): insert anonymized artifact URL.
\else
The artifact is public~\cite{artifact2026}.
\fi
It contains: every experiment script (\texttt{experiments/}), every result
JSON this paper cites (\texttt{experiments/results/}), the perturbation
tools (\texttt{tools/}), the machine-generated paper tables and their
generator (\texttt{scripts/build\_paper\_tables.py}), a truth ledger and
validity-boundary statement, per-file data-provenance documentation, and
both paper sources (this one and the long-form v9 draft, whose appendix
carries the full feature-manifest schema alongside
\texttt{spec/manifest-schema.json}).

\textbf{Claim verification.} \texttt{scripts/check\_claims.py} re-derives
every headline number in the repository README (the same numbers reported
here) from the committed result JSONs and fails on any drift; at this
writing it passes 31/31 checks. Result tables in this paper are generated
from the same JSONs, so table cells cannot silently diverge from the
evidence.

\textbf{Reproduction without restricted data.}
\texttt{scripts/reproduce\_public.sh} runs everything that needs no
restricted data (the synthetic S-curve sanity check and the M8
Byzantine/SPRT simulations) and runs the F3 analysis automatically if the
gated GT-HarmBench CSV is present, skipping it gracefully otherwise.

\textbf{Reproduction with PAN 2012.} The corpus is distributed by the PAN
organizers and is not redistributed in the artifact; the artifact documents
how to obtain it and where to place it. With it in place, the detection,
frontier, sweep, control, and FP-substrate experiments
(\texttt{exp\_trajectory\_lift.py}, \texttt{exp\_m3\_*.py},
\texttt{exp\_perturbation\_*.py}, \texttt{exp\_fp\_substrate.py}) reproduce
every real-data claim; each fails gracefully with setup instructions if the
data is absent. No script prints raw corpus text; sanity-check files verify
this posture per experiment. The included result JSONs let readers audit
every number without any restricted data. Appendix~\ref{app:artifact} gives
environment details, exact commands, and expected outputs.

% =====================================================================
\bibliographystyle{plain}
\bibliography{paper/references}

% =====================================================================
\appendix
\section{Artifact Appendix}
\label{app:artifact}

\textbf{Environment.} Python 3.10+ with \texttt{numpy},
\texttt{scikit-learn}, \texttt{matplotlib} (pinned in
\texttt{requirements.txt}); no GPU; a laptop-class machine suffices (the
heaviest experiments train linear models over $\sim$700k messages and run in
tens of minutes). LaTeX (TeX Live) to build the paper.

\textbf{Setup.} \texttt{pip install -r requirements.txt} at the repository
root.

\textbf{Reproduce without restricted data.}
\texttt{bash scripts/reproduce\_public.sh} runs the synthetic S-curve
check and the M8 simulations; it runs F3 if the gated GT-HarmBench CSV is
present at the documented path, else skips it with an explanation. Expected
outputs land in \texttt{experiments/results/} and
\texttt{validation/synthetic/results/}.

\textbf{Reproduce real-data claims.} Obtain PAN 2012 from the PAN organizers
(documented in \texttt{data/pan12/README.md}), place the training XML and
predator list under \texttt{data/pan12/train/}, then run each of the
following under \texttt{experiments/}:
\begin{itemize}\footnotesize
\item \texttt{exp\_trajectory\_lift.py}
\item \texttt{exp\_m3\_author\_split.py}, \texttt{exp\_m3\_frontier.py}
\item \texttt{exp\_perturbation\_sweep.py}
\item \texttt{exp\_perturbation\_sweep\_author\_disjoint.py}
\item \texttt{exp\_perturbation\_model\_controls.py}
\item \texttt{exp\_perturbation\_second\_family.py}
\item \texttt{exp\_fp\_substrate.py}
\end{itemize}
Each writes JSON + figures to \texttt{experiments/results/} and fails
gracefully with setup instructions if the data is absent.

\textbf{Claim-check procedure.} \texttt{python3 scripts/check\_claims.py}
(expects 31/31 OK); \texttt{python3 scripts/build\_paper\_tables.py}
regenerates every result table in \texttt{paper/tables/} from the JSONs.

\textbf{Known limitations.} PAN 2012 and GT-HarmBench cannot be
redistributed; exact figure rendering may vary with matplotlib version;
LinearSVC training is deterministic under the pinned seed (20260514) but
floating-point differences across BLAS builds can perturb the last decimal
of some metrics.

\end{document}
